\chapter{Hemmende Faktoren und Widerstände in Organisationen}
\label{ch:barriers}

Die proaktive Identifikation und Adressierung von Widerständen ist für den Erfolg einer \DO-Transformation ebenso entscheidend wie die Förderung von Erfolgsfaktoren.
Selbst die besten Initiativen können scheitern, wenn tief verwurzelte organisatorische und kulturelle Hemmnisse unbehandelt bleiben.
Dieses Kapitel analysiert die häufigsten Barrieren, die dem Wandel entgegenwirken.

\section{Traditionelle Hierarchien und Silodenken}
Das hartnäckigste Hindernis für \DO ist die traditionelle funktionale Organisation, die in starren Hierarchien und Abteilungssilos denkt.
Entwicklungs"=, Betriebs"=, Qualitätssicherungs"= und Sicherheitsteams arbeiten in getrennten Einheiten mit eigenen Zielen und Kennzahlen.
Diese Struktur behindert naturgemäß die für \DO essenzielle, abteilungsübergreifende Zusammenarbeit~\cite{faa23}.
Bürokratische Freigabeprozesse, langwierige Übergaben und mangelnde Transparenz sind direkte Konsequenzen dieses Silodenkens~\cite{kha20}.
% TODO Check source
% Die Ineffizienz solcher Strukturen lässt sich empirisch belegen:
% So zeigen die State of \DO Reports dramatische Leistungsunterschiede zwischen Unternehmen mit integrierten und silobasierten Arbeitsweisen.
% Enthaltene Daten verdeutlichen, dass die Defizite im sozialen Subsystem, wie in Kapitel 2 theoretisiert, direkt zu messbaren negativen Ergebnissen im technischen Betrieb führen.

% \begin{table}[htbp]
% 	\centering
% 	\begin{tabular}{p{0.40\linewidth} p{0.25\linewidth} p{0.25\linewidth}}
% 		\hline
% 		Leistungsmetrik              & Low Performer (silobasiert)   & Elite Performer (integriert)  \\
% 		\hline
% 		Change Failure Rate          & 46--60\%                      & 0--15\%                       \\
% 		Deployment Frequency         & Einmal pro Monat bis Halbjahr & On"=Demand (mehrmals täglich) \\
% 		Mean Time to Restore Service & 1 Woche bis 1 Monat           & $<$ 1 Tag                     \\
% 		\hline
% 	\end{tabular}
% 	\caption{Leistungsmetriken~\cite{dor20}}
% \end{table}

% Diese Daten sind die quantifizierbaren Symptome eines dysfunktionalen soziotechnischen Systems: Organisationen, die in Silos verharren, produzieren nicht nur langsamer, sondern auch mit signifikant höheren Fehlerraten und längeren Wiederherstellungszeiten.
% Das ist eine direkte Folge mangelnder Kollaboration, unzureichender Automatisierung und einer fehlenden \DO-Kultur.

\section{Fehlende gemeinsame Sprache}

Entwicklungs"= und Betriebsteams entstammen historisch unterschiedlichen Subkulturen innerhalb der IT\@.
Sie verwenden verschiedene Fachjargons, haben unterschiedliche Perspektiven auf Probleme und nutzen unterschiedliche Werkzeuge.
Diese sprachlichen und konzeptionellen Unterschiede führen oft zu Missverständnissen und ineffizienter Kommunikation.
\DO erfordert das aktive Überwinden dieser Barrieren und die Etablierung einer gemeinsamen Sprache und eines geteilten Verständnisses für den gesamten \foreignlanguage{american}{value stream}~\cite{cup16}.

\section{Kulturelle Spannungsfelder zwischen Stabilität und Innovation}
Ein fundamentaler Zielkonflikt prägt traditionelle IT"=Organisationen und stellt ein zentrales kulturelles Spannungsfeld dar:

Das Betriebsteam (Ops) ist auf maximale Stabilität, Zuverlässigkeit und Sicherheit der Produktivsysteme optimiert.
Sein Erfolg wird daran gemessen, Ausfallzeiten zu minimieren und Risiken zu vermeiden.
Veränderungen werden daher oft als potenzielle Bedrohung wahrgenommen.

Das Entwicklungsteam (Dev) ist auf schnelle Innovation und Veränderung ausgerichtet.
Sein Erfolg wird an der Geschwindigkeit gemessen, mit der neue Funktionen für den Kunden bereitgestellt werden können~\cite{ped21}.

Dieser inhärente Konflikt führt zu Reibungen und einer \enquote{Wir"=gegen"=die}"=Mentalität.
\DO versucht, genau dieses Spannungsfeld aufzulösen, indem es gemeinsame Ziele etabliert und beide Teams in eine gemeinsame Verantwortung nimmt.

\paragraph{}
Die Analyse dieser Hemmnisse zeigt, dass \DO weit mehr ist als die Einführung neuer Praktiken.
Es erfordert eine grundlegende Neuausrichtung organisationaler Strukturen und kultureller Normen, was im folgenden Kapitel eingehender diskutiert wird.